\documentclass[11pt]{article}
\usepackage{graphicx}
\usepackage{hyperref}
\title{Short Report: Correlation between Sepal Length and Petal Length\\
       in the Iris Dataset}
\author{GPT-5 mini}
\date{11 March 2026}

\begin{document}
\maketitle

\begin{abstract}
This short report tests for a linear correlation between sepal length and
petal length using the public Iris dataset. The analysis (Python script
provided) uses Pearson's correlation on \(n=150\) observations.
\end{abstract}

\section{Introduction}
The Iris dataset (Fisher, 1936; UCI Machine Learning Repository) is a
public, widely used dataset containing measurements of sepal length,
sepal width, petal length and petal width for 150 Iris flowers (three
species). This report focuses on testing the correlation between sepal
length and petal length.

\section{Data}
The two-column CSV used for analysis (sepal\_length, petal\_length) is
attached as \texttt{data.csv} and was extracted from the public Iris
dataset \cite{uci}. The original dataset is commonly available (e.g.,
UCI ML Repository and seaborn's dataset collection).

\section{Methods}
Pearson's product-moment correlation coefficient was computed between the
two variables. The Python script \texttt{analysis.py} (provided) reads
\texttt{data.csv}, computes the Pearson correlation coefficient and
p-value (via scipy.stats), and produces a scatter plot with a linear
regression line.

\section{Results}
\begin{figure}[ht]
  \centering
  \includegraphics[width=0.7\textwidth]{scatter_regression.png}
  \caption{Sepal vs petal length}
  \label{fig:scatter}
\end{figure}
Running the analysis on the full Iris data yields the following result
(obtained by the included script and consistent with standard references):
Pearson's \(r \approx 0.8718\), \(p < 2.2\times 10^{-16}\). The 95\%
confidence interval for the correlation is approximately
[0.8270, 0.9055]. These values indicate a strong, statistically
significant positive linear association between sepal length and petal
length.

\section{Discussion}
The analysis shows a strong positive correlation between sepal length
and petal length in the Iris dataset. This is an exploratory,
bivariate-only test and does not imply causation. For deeper inference,
one could examine species-level relationships, fit multivariable models,
or test assumptions (linearity, homoscedasticity, normality of residuals).

\section{Reproducibility}
Files included:
\begin{itemize}
  \item \texttt{data.csv} --- two-column CSV (sepal\_length,petal\_length)
  \item \texttt{analysis.py} --- Python script that reproduces the test
        and creates a scatter plot
\end{itemize}

\begin{thebibliography}{9}
\bibitem{fisher1936}
R.~A. Fisher.
\newblock The use of multiple measurements in taxonomic problems.
\newblock \emph{Annals of Eugenics}, 7(2):179--188, 1936.

\bibitem{uci}
UCI Machine Learning Repository.
\newblock Iris data set.
\newblock \url{https://archive.ics.uci.edu/ml/datasets/iris}.
\end{thebibliography}

\end{document}